\documentclass[12pt]{article}
\usepackage[utf8]{inputenc}
\usepackage[T1]{fontenc}
\usepackage[margin=1in]{geometry}
\usepackage{setspace}
\usepackage{xcolor}
\usepackage[colorlinks=true, urlcolor=blue!60!black, linkcolor=blue!60!black]{hyperref}
\setlength{\parindent}{0em}
\setlength{\parskip}{0.6em}
\renewcommand{\baselinestretch}{1.0}
\begin{document}
\begin{center}
{\LARGE \textit{New York Times Co. v. Sullivan}, 376 U.S. 254 (1964)\par}
\vspace{0.5em}
{\large Case Brief --- Two-Page Edition\par}
\end{center}

\section{Memory Jogger}
A public official may recover for a defamatory falsehood about his official conduct only by proving ``actual malice''---knowing or reckless falsity (\href{https://supreme.justia.com/cases/federal/us/376/254/}{Justia}).

\section{Detailed Case Facts}
In 1960 the \textit{New York Times} ran a civil-rights fundraising advertisement, ``Heed Their Rising Voices,'' that contained minor errors about events in Montgomery, Alabama. City Commissioner L.B. Sullivan, who supervised the police but was not named, sued the \textit{Times} and four clergymen listed as endorsers.

\section{Procedural History}
Instructed that the statements were libelous per se, so that malice and damages were presumed, the jury awarded \$500,000. The Supreme Court of Alabama affirmed.

\section{Judicial Votes}
Reversed, 9--0. Brennan wrote for the Court; Black (with Douglas) concurred, and Goldberg (with Douglas) concurred in the result.

\section{Holding}
A public official must prove with convincing clarity that a defamatory falsehood about his official conduct was made with knowledge of its falsity or reckless disregard of the truth. The evidence here showed neither actual malice nor that the ad was ``of and concerning'' Sullivan.

\section{Analysis of Opinions}
\textbf{Majority:} Libel has no ``talismanic immunity.'' Errors are inevitable in free debate and need ``breathing space''; the Sedition Act controversy defines the First Amendment's central meaning, and huge civil awards chill criticism more than criminal sedition laws. A truth defense alone induces self-censorship. Failing to check the paper's own files was at most negligence, and treating criticism of government as libel of officials revived seditious libel. \textbf{Black and Goldberg:} criticism of official conduct deserves absolute protection, since actual malice is elusive and juries are hostile. \emph{Later:} \textit{Curtis Publishing} and \textit{Gertz} extended and bounded the rule.

\section{Examples: Future Applications}
\textbf{Same-Side:} an op-ed with careless errors about an Attorney General is not actionable. \textbf{Opposite-Side:} a blogger who purposefully avoids records refuting a charge against a police chief acts with reckless disregard. \textbf{Fence-Sitter:} a business owner who testified once at a zoning hearing and is falsely accused of bribery.

\section{Critique}
Originalists argue that the Amendments never displaced state libel law; absolutists say the actual-malice inquiry still chills speech. Pragmatists worry it under-deters deliberate falsehoods online.

\section{Key Quotations}
\begin{quote}
``The constitutional guarantees require, we think, a federal rule that prohibits a public official from recovering damages for a defamatory falsehood relating to his official conduct unless he proves that the statement was made with `actual malice'---that is, with knowledge that it was false or with reckless disregard of whether it was false or not.''
\end{quote}
\noindent --- Justice Brennan, for the Court, 376 U.S. at 279--80.

\end{document}
